\scanpage{100}
\mspage{49}
Prouvons (si nous pouvons) ce complément \unsure{sympa}. Je précise que j'ai eu à faire l'hypothèse de $W$-prélissité \inserted{sur $X/S$} (ou de $W$-prépropreté), pour être
sûr que l'image inverse $X\times_S S'$, par $f:S'\unsure{\relbar|\longrightarrow}S$, coïncide bien avec $f^*$ pour les $\mathrm{HOT}_W^{\mathrm{lc}}$ -- autrement je n'ai aucune chance. ``On sait'' (tout est à revoir !) que c'est OK
dans les trois cas $X$ $W$-lisse, $X$ $W$-propre, $X$ une $W$-fibration (mais en quel sens ???). \unsure{La} vérification à faire \inserted{(quand on travaille avec les $W_S^d$)} est celle\unsure{-| }ci :
\[
X\overset{i_X}{\longrightarrow}\Phi_S(X)=Y\ \text{est dans }W_S^d\ \text{(sans mérite), et}\tag*{(*)}
\]
le reste par tout changement de base.

Mais pour des $X,Y\in\mathrm{Ob}\,\mathrm{Cat}_{/S}$ qui sont $W$-pré-lisses, les $W_S^d$-équivalences coïncident avec $W_S^f$, laquelle est préservée par chgt de base, ainsi que l'hyp. sur $X,Y$. Donc on est tranquille.

On est ramené à prouver que si $X$ est $W$-prélisse \inserted{sur $S$} \deleted{et si} et une $W$-pseudofibration, \deleted{alors} et si $S'\longrightarrow S$ est un \unsure{homotopisme}, alors $X'\longrightarrow X$ est dans $W$. Par l'argument de \deleted{\unsure{tout}} tantôt utilisant $X\longrightarrow\Phi_S(X)$, on se ramène au cas où $X$ est fibrée scindée
\deleted{Prenant \unsure{une|un} inverse homotopique $g$ de $f$, considérons $g^*(X')=(\underbrace{fg}_{\varphi})^*(X)=X_{\unsure{1|1'}}$}
\deleted{\diagram{\begin{tikzcd}[ampersand replacement=\&]
S'\arrow[r,"{\unsure{f}}"] \& S\arrow[r,"g"] \& S'\arrow[r,"f"] \& S,
\end{tikzcd}}}
% archive sheet number: 100
